\label{ch:detector-implementation}

This chapter develops every formula \texttt{tpeanuts.detector.sno\_ii}
implements: the 38-observable structure (\cref{sec:38-observables}), the
equivalent-flux definition each channel shares
(\cref{sec:equivalent-flux}), the real salt-phase energy response
(\cref{sec:response-function}), and the day/night composition with the
Earth-matter regeneration machinery shared by the rest of \texttt{tpeanuts}
(\cref{sec:daynight-composition}).

\section{The 38-observable structure}
\label{sec:38-observables}

\modulehead{detector/sno\_ii/parameters.py}{\texttt{CHANNEL\_ORDER},
binning, and every other real constant this chapter's formulas need.}

\begin{theorybox}
SNO's own unconstrained (model-independent) salt-phase signal extraction
produces, separately for day and night, a 17-bin CC recoil-electron
spectrum (16 bins of 0.5~MeV from 5.5-13.5~MeV, one wide bin
13.5-20.0~MeV) plus single integrated NC and ES equivalent fluxes -- 19
observables per period, 38 total:
\begin{equation}
  \keyresult{\vec Y = \bigl(Y_{\rm NC}, Y_{{\rm CC},1},\dots,Y_{{\rm CC},17}, Y_{\rm ES}\bigr)_{\rm day}
    \;\oplus\;
    \bigl(Y_{\rm NC}, Y_{{\rm CC},1},\dots,Y_{{\rm CC},17}, Y_{\rm ES}\bigr)_{\rm night}.}
  \label{eq:observable-vector}
\end{equation}
Every real published table this document's data was transcribed from uses
this same ordering (\cref{ch:data-release}) except Table XXXIV specifically,
which prints its own columns as (NC, ES, CC$_1$,\dots,CC$_{17}$) -- a
genuine, disclosed difference from \labelcref{eq:observable-vector}'s own
order, kept as printed (\cref{sec:covariance-construction}) rather than
silently reordered, to keep future spot-checks against the primary source
straightforward.
\end{theorybox}

\begin{implbox}
\code{CHANNEL\_ORDER = ("NC", "CC1", ..., "CC17", "ES")} (19 entries);
\code{N\_OBSERVABLES\_TOTAL = 38}; \code{CC\_BIN\_EDGES\_MEV} holds the 17
bin edges above. \pyfunc{load\_observed\_vector\_and\_covariance} (Chapter 4)
assembles the real day-then-night \labelcref{eq:observable-vector} vector
directly from the tables of \cref{ch:data-release}.
\end{implbox}

\section{The equivalent-flux definition}
\label{sec:equivalent-flux}

\modulehead{detector/sno\_ii/observable.py}{\pyfunc{cc\_equivalent\_flux\_spectrum},
\pyfunc{es\_equivalent\_flux}, \pyfunc{nc\_equivalent\_flux}.}

\begin{theorybox}
Every observable in \labelcref{eq:observable-vector} is an \emph{equivalent
$^8$B flux}, not a raw event count (Eq.~A1 of \cite{Aharmim2005SaltPhase}):
for the CC channel's bin $i$,
\begin{equation}
  \keyresult{Y_{{\rm CC},i} = \Phi_{\rm tot}\;
    \frac{\displaystyle\int dE_\nu\int dT_e\int_{{\rm bin}\,i} dT_{\rm eff}\;
      \phi(E_\nu)\,P_{ee}(E_\nu)\,\frac{d\sigma_{\rm CC}}{dT_e}(E_\nu,T_e)\,R(T_e,T_{\rm eff})}
    {\displaystyle\int dE_\nu\int dT_e\int_{5.5}^{\infty} dT_{\rm eff}\;
      \phi(E_\nu)\,\frac{d\sigma_{\rm CC}}{dT_e}(E_\nu,T_e)\,R(T_e,T_{\rm eff})},}
  \label{eq:cc-equivalent-flux}
\end{equation}
with $\Phi_{\rm tot}=\int\phi(E_\nu)\,dE_\nu$ (the total injected flux this
call's $\phi(E_\nu)$ integrates to) and $\phi(E_\nu)$ the total (not yet
survival-probability-weighted) differential $^8$B(+hep) flux. Physically:
the \emph{denominator} is the same detector fold evaluated at
$P_{ee}\equiv1$ (an undistorted, unoscillated spectrum) and integrated over
the \emph{entire} analysis range above the 5.5~MeV threshold, so the ratio
answers ``what fraction of the total unoscillated rate would this bin's
oscillated rate correspond to''. A direct, exact consequence: summing
\labelcref{eq:cc-equivalent-flux} over every bin at $P_{ee}\equiv1$ recovers
$\Phi_{\rm tot}$ exactly, since the numerator's sum over all bins becomes
identical to the denominator's own full-range integral.

ES mirrors \labelcref{eq:cc-equivalent-flux} with the elastic-scattering
cross sections and $(1-P_{ee})$ weighting the $\nu_{\mu,\tau}$ contribution
in the numerator, while the denominator stays a pure-$\nu_e$ ($P_{ee}\equiv1$)
reference -- SNO's own ``equivalent electron-neutrino flux'' convention for
this channel:
\begin{equation}
  Y_{\rm ES} = \Phi_{\rm tot}\;
    \frac{\int dE_\nu\int dT_e\int_{5.5}^{\infty} dT_{\rm eff}\;\phi(E_\nu)
      \bigl[P_{ee}\frac{d\sigma_e}{dT_e}+(1-P_{ee})\frac{d\sigma_x}{dT_e}\bigr]R(T_e,T_{\rm eff})}
    {\int dE_\nu\int dT_e\int_{5.5}^{\infty} dT_{\rm eff}\;\phi(E_\nu)\,
      \frac{d\sigma_e}{dT_e}(E_\nu,T_e)\,R(T_e,T_{\rm eff})}.
  \label{eq:es-equivalent-flux}
\end{equation}

NC needs no detector fold at all (Eq.~10 of \cite{Aharmim2005SaltPhase}):
SNO's own signal extraction already un-folds the neutron-capture response
for this channel, so
\begin{equation}
  \keyresult{Y_{\rm NC} = \int dE_\nu\;\phi(E_\nu)\sum_{\alpha\,\in\,{\rm active}} P_{e\alpha}(E_\nu)
    = \Phi_{\rm tot}\quad\text{(Standard Model, three active flavours).}}
  \label{eq:nc-equivalent-flux}
\end{equation}
\end{theorybox}

\begin{implbox}
\pyfunc{cc\_equivalent\_flux\_spectrum(probabilities, flux\_tot\_MeV,
bin\_edges\_MeV, ...)} implements \labelcref{eq:cc-equivalent-flux} by
reusing \texttt{detector.common.event\_rate.predicted\_counts} for
\emph{both} the numerator (per bin) and the denominator (one wide bin, full
range), with \code{n\_target=1}, \code{exposure=1} -- since both the
numerator and denominator share the identical (arbitrary) normalization
convention, it cancels exactly in the ratio, so no dedicated
``un-normalized fold'' helper is needed. \pyfunc{es\_equivalent\_flux}
mirrors this with \texttt{nue\_cross\_section\_grid}/
\texttt{numutau\_cross\_section\_grid}. \pyfunc{nc\_equivalent\_flux(probabilities,
flux\_tot\_MeV, E\_nu\_grid\_MeV)} implements \labelcref{eq:nc-equivalent-flux}
as a single \texttt{torch.trapezoid} call, with an optional
\code{probabilities=None} fast path assuming unit active probability
(future sterile-neutrino extension point, mirroring
\texttt{detector.sno.event\_rate.nc\_event\_rate}'s own convention).
\end{implbox}

\begin{resultbox}
Three exact closure properties, verified in
\texttt{detector/sno\_ii/test/test2\_observable.py} against the real
$^8$B flux (not synthetic inputs): at $P_{ee}\equiv1$, the CC spectrum's
own sum recovers $\Phi_{\rm tot}$ to within $10^{-6}$ relative; at an
energy-independent $P_{ee}=p$, both the CC spectrum and the ES flux scale
\emph{exactly} linearly in $p$ (not merely approximately, a direct
consequence of \labelcref{eq:cc-equivalent-flux,eq:es-equivalent-flux}
being linear folds of a linear input); the NC flux's gradient with respect
to $P_{ee}$ is exactly zero for any energy-dependent $P_{ee}(E)$ shape that
keeps the three active probabilities summing to 1, confirming
\labelcref{eq:nc-equivalent-flux}'s oscillation-independence numerically,
not just analytically.
\end{resultbox}

\section{The salt-phase energy response}
\label{sec:response-function}

\modulehead{detector/sno\_ii/response.py}{\pyfunc{sigma\_MeV},
\pyfunc{response\_matrix}.}

\begin{theorybox}
The Gaussian response function in \labelcref{eq:cc-equivalent-flux,eq:es-equivalent-flux},
\begin{equation}
  R(T_e,T_{\rm eff}) = \frac{1}{\sqrt{2\pi}\sigma_T}
    \exp\!\left[-\frac{(T_e-T_{\rm eff})^2}{2\sigma_T^2}\right],
  \label{eq:response-gaussian}
\end{equation}
uses the salt-phase's own real kinetic-energy resolution parametrisation
(Eq.~A3 of \cite{Aharmim2005SaltPhase}, verified directly against the
primary source, distinct from -- and not to be confused with -- both
Phase~I's own coefficients and the earlier flux-only salt paper's
methodology note \parencite{Ahmed2004SaltFlux}, which quotes a visibly
different formula for a different, non-spectral extraction):
\begin{equation}
  \keyresult{\sigma_T(T_e) = -0.131 + 0.383\sqrt{T_e} + 0.03731\,T_e \quad[{\rm MeV},\ T_e\text{ in MeV}].}
  \label{eq:response-sigma}
\end{equation}
\end{theorybox}

\begin{implbox}
\pyfunc{response\_matrix(T\_grid\_MeV, Tprime\_grid\_MeV, ...)} reuses
\texttt{detector.common.response.gaussian\_response\_matrix} unchanged
(no new generic response machinery), evaluated with
\labelcref{eq:response-sigma}'s coefficients on a shared 500-point, 0-20~MeV
grid. Verified: $\sigma_T(10\,{\rm MeV})=1.453$~MeV (consistent with
Phase~I's own commonly-quoted $\approx1.40$~MeV at the same energy), and
the response integrates to $0.999999999994$ over reconstructed energy at
an interior grid point -- a genuine probability-conserving kernel, not an
approximation with a visible normalization leak.
\end{implbox}

\section{Day/night composition and Earth-matter regeneration}
\label{sec:daynight-composition}

\modulehead{detector/sno\_ii/inference\_model.py}{\pyclass{SNOPhaseIIObservableModel}.}

\begin{theorybox}
Day and night flavour probabilities are built exactly as for this project's
own Phase-I \pyclass{SNODayNightModel}: MSW production-region mixing inside
the Sun \parencite{MikheyevSmirnov1985} followed by Earth-matter
regeneration \parencite{Wolfenstein1978}, averaged over SNO's own real
60-bin $\cos(\text{zenith})$ livetime distribution (Table XXXI) split at
the horizon,
\begin{equation}
  \keyresult{P_{ee}^{\rm day,night}(E) = \sum_\eta w_\eta^{\rm day,night}\,
    P_{ee}\bigl(E,\eta;\theta_{12},\Delta m^2_{21}\bigr),}
  \label{eq:daynight-average}
\end{equation}
with $w_\eta$ the real, re-normalized (summing to 1 per period) exposure
weight at nadir angle $\eta$. Since $^8$B and hep are produced at very
slightly different radii inside the Sun,
\texttt{solar\_probability\_mass} formally depends on the production
source -- this model evaluates \labelcref{eq:daynight-average} once using
$^8$B's own production profile and applies the resulting $P_{ee}(E)$ to
the \emph{combined} $^8$B+hep flux, an explicit, quantified-as-negligible
approximation (hep contributes $9.3\times10^3\,{\rm cm^{-2}s^{-1}}$
against $^8$B's $\sim4$--$5\times10^6\,{\rm cm^{-2}s^{-1}}$, $\sim0.15\%$
of the total).
\end{theorybox}

\begin{notebox}
Combining the two sources' \emph{fluxes} before folding through
\labelcref{eq:cc-equivalent-flux,eq:es-equivalent-flux} (rather than
computing two separate equivalent-flux predictions and adding the results)
is required for correctness, not merely convenient: since
\labelcref{eq:cc-equivalent-flux} is a ratio,
$(A_1+A_2)/(B_1+B_2)\neq A_1/B_1+A_2/B_2$ in general -- summing separately
computed per-source equivalent fluxes would silently give the wrong answer.
\end{notebox}

\section{Free \texorpdfstring{$^8$B}{8B} normalization}
\label{sec:free-normalization}

\begin{theorybox}
Mirroring SNO's own free $f_B$ (Section~XIII of
\cite{Aharmim2005SaltPhase}: ``the $^8$B flux was free in the fit''), this
model exposes
\begin{equation}
  \keyresult{\Phi_{8B} = \Phi_{8B}^{\rm ref}\,e^{\lambda_{8B}},}
  \label{eq:free-normalization}
\end{equation}
with $\Phi_{8B}^{\rm ref}$ the Standard Solar Model table value already
loaded by \texttt{source.solar.SolarNeutrinoSource}, so $\lambda_{8B}=0$
reproduces the SSM prediction exactly and the fitted flux stays strictly
positive by construction, avoiding the need for a separate positivity
constraint in the optimizer.
\end{theorybox}

\begin{implbox}
\pyclass{SNOPhaseIIObservableModel.free} $=$
\code{oscillation\_model.free + ("log\_phi\_8B",)} -- three free parameters
total, $(\theta_{12},\Delta m^2_{21},\lambda_{8B})$. \code{.predict(theta)}
builds \labelcref{eq:free-normalization}, combines the $^8$B+hep flux, calls
\labelcref{eq:daynight-average}, and returns the concatenated 38-observable
\labelcref{eq:observable-vector}, all in units of $10^6\,{\rm cm^{-2}s^{-1}}$
(\texttt{FLUX\_UNIT\_CM2S}), matching the tabulated data's own units --
\texttt{SolarNeutrinoSource.total\_flux} itself returns a raw (not
$10^6$-scaled) cm$^{-2}$s$^{-1}$ value, so this conversion is applied
explicitly rather than left implicit.
\end{implbox}
